% 从 resume-zh_CN.tex 移出的项目底稿（2026-10-31），要恢复时整段粘回「项目经历」下即可

\datedsubsection{\textbf{Springboot服务单机启动加速}}{技术负责人}
\begin{onehalfspacing}
\textbf{项目背景}：服务的依赖越多启动速度越慢，当业务系统随着迭代的进行依赖越来越多启动速度会越来越慢，极大影响开发测试时的部署速度从而降低迭代效率，同时在线上出事故需要回滚的时候，服务启动时间越长，事故时间越长影响越大，所以提升服务单机启动速度对于开发和生产都有着重大意义。\\

\textbf{工作内容}：
\begin{itemize}
  \item 梳理启动流程：启动JVM -> 类加载 -> 创建ApplicationContext -> Bean初始化 -> 启动内置Web服务器。
  \item 定位性能瓶颈：通过在启动流程各阶段打点上报耗时的方式，定位到类加载这一步骤占整个启动流程的大部分。
  \item 分析瓶颈原因：JVM在URLClassLoader类加载过程中会调用findClass方法查找类属于哪个jar包，原生实现的查找方式会依次在所有jar包中查找目标类，是线性时间复杂度，对于依赖较多的服务而言这部分耗时长。
  \item 制定解决方案：借鉴用空间换时间的思想，我们在类加载之前首先构建JarIndex(ClassName -> JarPath)，然后在类加载时直接根据构建的JarIndex查找定位目标类(用反射替换URLClassLoader中的核心数据结构UCP)，将时间复杂度由O(N)降低为O(1)，同时我们提供了多种不同类型的索引结构供不同服务使用（Map, TrieTree, 多级Map）。
  \item 打包解决方案：将服务启动加速的解决方案打包供不同业务服务即插即用，同时提供配置化参数的能力让不同业务服务能根据自己需求选用不同的JarIndex数据类型。\\
\end{itemize}

\textbf{项目成果}：
\begin{itemize}
  \item 启动加速效果：经过测试，在引入服务启动加速方案后，服务启动的平均耗时减小30\%到50\%（与服务所加载类的多少有关，服务加载类越多启动加速越明显）。
  \item 方案推广应用：服务加速启动方案通过打包的方式成功推广到部门各个团队的～50个服务，都取得了不错的效果，获得广泛好评。\\
\end{itemize}
\end{onehalfspacing}

